\documentclass[11pt,a4paper]{article}
\usepackage{amsmath,amsthm,amssymb,mathtools}
\usepackage[margin=2.5cm]{geometry}
\usepackage{hyperref}
\usepackage{enumitem}
\usepackage{tikz}
\usepackage{xcolor}

\hypersetup{
  colorlinks=true,
  linkcolor=blue!70!black,
  citecolor=green!50!black,
  urlcolor=blue!60!black
}

% Theorem environments
\newtheorem{theorem}{Theorem}[section]
\newtheorem{lemma}[theorem]{Lemma}
\newtheorem{proposition}[theorem]{Proposition}
\newtheorem{corollary}[theorem]{Corollary}
\theoremstyle{definition}
\newtheorem{definition}[theorem]{Definition}
\newtheorem{example}[theorem]{Example}
\newtheorem{problem}{Open Problem}
\theoremstyle{remark}
\newtheorem{remark}[theorem]{Remark}

\newcommand{\N}{\mathbb{N}}
\newcommand{\Q}{\mathbb{Q}}
\newcommand{\R}{\mathbb{R}}
\newcommand{\PA}{\mathsf{PA}}
\newcommand{\ZFC}{\mathsf{ZFC}}
\newcommand{\Prov}{\mathrm{Prov}}
\newcommand{\Bel}{\mathrm{Bel}}
\newcommand{\EU}{\mathrm{EU}}
\newcommand{\Do}{\mathrm{do}}
\newcommand{\Tr}{\mathrm{Tr}}

\title{\textbf{Embedded Agency and Formal Decision Theory:\\
Mathematical Foundations for Agentic AI}\\ [0.5em]
\large Research Notes --- ETH Zurich ML Lecture Series}
\author{Agentic AI Frontier Seminar}
\date{Spring 2026}

\begin{document}
\maketitle

\begin{abstract}
These notes survey the mathematical theory of \emph{embedded agency}: the
study of agents that reason about environments of which they are
themselves a part.  We develop the core formalisms---logical induction,
functional decision theory, L\"ob's theorem and the L\"obian obstacle,
reflective oracles, and Cartesian frames---with full definitions and
proof sketches aimed at a mathematically mature audience.  We then
connect these foundations to contemporary agentic AI systems (tool-using
LLMs, multi-agent architectures, constitutional AI) and close with open
problems at the research frontier.
\end{abstract}

\tableofcontents
\newpage

%==========================================================================
\section{The Dualism Problem}\label{sec:dualism}
%==========================================================================

Classical AI and decision theory rest on a \emph{Cartesian} split:
the agent is \emph{outside} the environment, receiving observations and
emitting actions through a clean interface.  Formally, one posits a
function $\pi : \mathcal{O}^* \to \mathcal{A}$ (the policy), a
transition kernel $T : \mathcal{S} \times \mathcal{A} \to
\Delta(\mathcal{S})$, and an observation function $O : \mathcal{S} \to
\mathcal{O}$, where crucially the agent's internal state is \emph{not}
an element of $\mathcal{S}$.

\begin{definition}[Dualistic agent model]
A \emph{dualistic agent--environment pair} is a tuple
$(\mathcal{S}, \mathcal{A}, \mathcal{O}, T, O, \pi)$ where the
agent $\pi$ and the environment $(S, T, O)$ share no state.
The agent can be ``lifted out'' and replaced without altering
the environment dynamics.
\end{definition}

This separation breaks down the moment we acknowledge four facts:
\begin{enumerate}[label=(\roman*)]
  \item \textbf{Physical embedding.} The agent is made of atoms in the
    environment.  Its hardware can be damaged, copied, or inspected.
  \item \textbf{Self-reference.} The agent may reason about its own
    source code or weights, creating G\"odelian fixed-point phenomena.
  \item \textbf{Logical uncertainty.} Even with full knowledge of the
    environment's laws, the agent faces irreducible uncertainty about
    logical and mathematical consequences.
  \item \textbf{Subsystem alignment.} The agent's goals may conflict
    with goals of sub-agents it creates or of the larger system it
    inhabits.
\end{enumerate}

Demski and Garrabrant~\cite{demski2019embedded} organise the embedded
agency research programme around four corresponding sub-problems:
\begin{itemize}
  \item \textbf{Decision problem:} How should an embedded agent choose
    actions when it cannot separate ``intervention'' from ``observation''?
  \item \textbf{Logical uncertainty problem:} How should an agent assign
    credences to logical claims it cannot yet verify?
  \item \textbf{High-bandwidth problem:} The agent is a \emph{small}
    part of a \emph{large} environment---it cannot simulate the whole.
  \item \textbf{Subsystem alignment problem:} How can an agent trust
    future versions of itself, or sub-agents it spawns?
\end{itemize}

The remainder of these notes addresses each sub-problem through a
concrete mathematical framework.

%==========================================================================
\section{Logical Induction}\label{sec:logical-induction}
%==========================================================================

\subsection{Motivation: beyond Bayesian epistemology}

A Bayesian agent assigns a prior over possible worlds and updates by
conditioning.  But for an embedded agent, many questions are
\emph{logical} rather than empirical: ``Is $P \neq NP$?'', ``Does my
own source code halt on input 42?''  A coherent prior over logical
sentences requires logical omniscience (assigning probability~1 to all
theorems), which is computationally impossible.

Logical induction, introduced by Garrabrant et al.~\cite{garrabrant2016logical},
provides a computable process whose beliefs about logical sentences
converge to satisfy a rich collection of rationality desiderata.

\subsection{The trading framework}

Fix a first-order language $\mathcal{L}$ (e.g.\ the language of Peano
arithmetic $\PA$) and an enumeration $\phi_1, \phi_2, \ldots$ of its
sentences.

\begin{definition}[Pricing]
A \emph{pricing} is a function $\mathbb{P} : \mathcal{L} \to [0,1]$.
We interpret $\mathbb{P}(\phi)$ as the ``price of a share'' that pays
\$1 if $\phi$ is true in the standard model.
\end{definition}

\begin{definition}[Trader]
A \emph{trader} is a computable function
$\mathcal{T} : [0,1]^{<\omega} \to \mathcal{L} \times \R$
that, given the history of past prices, outputs a sentence $\phi$
and a stake $\alpha$ (positive = buy, negative = sell).
The \emph{value} of trader $\mathcal{T}$'s position after day $n$
relative to a world (a consistent completion $W$ of the deductive
system) is:
\[
  V_n^W(\mathcal{T}) \;=\; \sum_{i=1}^{n} \alpha_i \bigl(
    \mathbf{1}_W[\phi_i] - \mathbb{P}_i(\phi_i)
  \bigr),
\]
where $\mathbf{1}_W[\phi_i] = 1$ if $\phi_i \in W$ and $0$ otherwise.
\end{definition}

\begin{definition}[Exploiting a market]
A trader $\mathcal{T}$ \emph{exploits} a sequence of pricings
$(\mathbb{P}_n)_{n \geq 1}$ if there exists a constant $c$ such that
for every consistent completion $W$:
\[
  V_n^W(\mathcal{T}) \;\geq\; -c \quad\text{for all } n,
  \qquad\text{and}\qquad
  \limsup_{n \to \infty}\, V_n^W(\mathcal{T}) \;=\; +\infty.
\]
That is, $\mathcal{T}$ has bounded downside and unbounded upside.
\end{definition}

\begin{definition}[Logical induction criterion]\label{def:lic}
A computable sequence of pricings $(\mathbb{P}_n)_{n \geq 1}$ is a
\emph{logical inductor} over a deductive system $D$ if no
polynomially-computable trader exploits it.
\end{definition}

The polynomial-time bound on traders is essential: it models the
resource limitations of an embedded reasoner while still capturing a
vast class of ``efficiently noticeable'' patterns.

\subsection{Existence and main properties}

\begin{theorem}[Garrabrant et al.\ 2016]\label{thm:li-existence}
Logical inductors exist.  Specifically, there is a computable sequence
of pricings $(\mathbb{P}_n)_{n \geq 1}$ satisfying
Definition~\ref{def:lic}.
\end{theorem}

\begin{proof}[Proof sketch]
Enumerate all poly-time traders $\mathcal{T}_1, \mathcal{T}_2, \ldots$.
On day $n$, let trader $\mathcal{T}_k$ (for $k \leq n$) submit demands.
Set prices by solving a linear programme (or a fixed-point equation)
that minimises the aggregate ``profit opportunity'' of all traders seen
so far, subject to prices lying in $[0,1]$.  Because the weight of each
trader shrinks (e.g.\ weight $2^{-k}$ for $\mathcal{T}_k$), the total
exploitation is bounded above, and no single trader can achieve
unbounded profit.
\end{proof}

\begin{theorem}[Properties of logical inductors]\label{thm:li-props}
Let $(\mathbb{P}_n)$ be a logical inductor.  Then:
\begin{enumerate}[label=(\alph*)]
  \item \textbf{Convergence.} For every sentence $\phi$,
    $\lim_{n \to \infty} \mathbb{P}_n(\phi)$ exists.  Write
    $\mathbb{P}_\infty(\phi)$ for the limit.
  \item \textbf{Coherence.} $\mathbb{P}_\infty$ is a finitely additive
    probability measure on $\mathcal{L}$ (modulo provable equivalence).
  \item \textbf{Calibration.} The prices are asymptotically calibrated:
    among sentences assigned limiting price $p$, the fraction that are
    true converges to $p$.
  \item \textbf{Gaifman condition.} If $D \vdash \phi$, then
    $\mathbb{P}_\infty(\phi) = 1$.  (No probability mass on refutable
    sentences.)
  \item \textbf{Non-dogmatism.} If $\phi$ is consistent with $D$, then
    $\mathbb{P}_\infty(\phi) > 0$.
  \item \textbf{Domination of computably enumerable sequences.}
    If a CE sequence of pricings $(\mathbb{Q}_n)$ is itself eventually
    coherent, then $\mathbb{P}$ eventually assigns at least as much
    probability to true sentences as $\mathbb{Q}$ does.
\end{enumerate}
\end{theorem}

\begin{remark}
Property (f) is the logical analogue of Blackwell--Dubins merging: the
logical inductor eventually dominates any ``reasonable'' computable
belief sequence.  This is the sense in which logical induction is
\emph{optimal} under computational constraints.
\end{remark}

\subsection{Relationship to prediction markets}

Logical induction can be understood as an internalised prediction market
where each trader represents a ``pattern of reasoning.'' The market
price aggregates these patterns efficiently.  The polynomial-time
restriction models bounded rationality, and the no-exploitation
condition is a computationally bounded analogue of the \emph{efficient
market hypothesis}.

%==========================================================================
\section{Functional Decision Theory}\label{sec:fdt}
%==========================================================================

\subsection{Three decision theories}

We consider an agent choosing an action $a \in \mathcal{A}$ in a
decision problem with outcome function
$\text{Out}: \mathcal{A} \times \mathcal{E} \to \mathcal{O}$, where
$\mathcal{E}$ is a set of environment states.

\begin{definition}[Causal Decision Theory (CDT)]
CDT recommends:
\[
  a^*_{\text{CDT}} = \arg\max_{a \in \mathcal{A}} \sum_{e \in \mathcal{E}}
    P(e \mid \Do(a))\, U(\text{Out}(a, e)),
\]
where $P(e \mid \Do(a))$ denotes the probability of $e$ under a causal
intervention setting the agent's action to $a$, using Pearl's
do-calculus.
\end{definition}

\begin{definition}[Evidential Decision Theory (EDT)]
EDT recommends:
\[
  a^*_{\text{EDT}} = \arg\max_{a \in \mathcal{A}} \sum_{e \in \mathcal{E}}
    P(e \mid a)\, U(\text{Out}(a, e)),
\]
where $P(e \mid a)$ is ordinary conditional probability.
\end{definition}

\begin{definition}[Functional Decision Theory (FDT)]
FDT (Yudkowsky and Soares~\cite{soares2018fdt}) recommends:
\[
  a^*_{\text{FDT}} = \arg\max_{a \in \mathcal{A}} \sum_{e \in \mathcal{E}}
    P(e \mid \widehat{\text{FDT}}() = a)\, U(\text{Out}(a, e)),
\]
where $\widehat{\text{FDT}}()$ is a rigid designator for the output of
the agent's \emph{decision algorithm}, and the conditional is over
\emph{logical} or \emph{subjunctive} dependence: ``what would the world
look like if the computation FDT() returned $a$?''
\end{definition}

The key distinction: CDT conditions on \emph{causal} intervention, EDT
on \emph{evidential} correlation, and FDT on \emph{logical/functional}
dependence of the environment on the agent's algorithm.

\subsection{Newcomb's Problem}

\begin{example}[Newcomb's Problem]
A predictor Omega, who has never been wrong, places \$1{,}000 in a
transparent box and either \$1{,}000{,}000 or \$0 in an opaque box,
depending on whether it predicted you would take one box or two.

\medskip
\begin{center}
\begin{tabular}{l|cc}
  & Omega predicted 1-box & Omega predicted 2-box \\
\hline
Take 1 box & \$1{,}000{,}000 & \$0 \\
Take 2 boxes & \$1{,}001{,}000 & \$1{,}000 \\
\end{tabular}
\end{center}
\medskip

\noindent
CDT two-boxes (causal dominance), EDT one-boxes (evidential correlation),
and FDT one-boxes (the computation FDT() returning ``one-box'' is the
same computation Omega simulated, so $P(\text{predicted 1-box} \mid
\widehat{\text{FDT}}() = \text{1-box}) \approx 1$).
\end{example}

\subsection{Prisoner's Dilemma with a copy}

\begin{example}
You play Prisoner's Dilemma against an exact copy of yourself (same
source code, same inputs).

CDT: defects (causal independence from the copy).

EDT: cooperates (my action is evidence of my copy's action).

FDT: cooperates (the same computation determines both outputs; if
$\widehat{\text{FDT}}() = C$, then both cooperate, yielding $(C,C)$).

\medskip
FDT achieves the cooperative payoff because it recognises that its
algorithm \emph{is} the same function being evaluated in both positions.
\end{example}

\subsection{Parfit's Hitchhiker}

\begin{example}
You are stranded in the desert.  A driver will rescue you if and only if
she predicts you will pay \$1{,}000 upon reaching town.  Once in town,
paying is costly and you have no further interaction with the driver.

CDT: does not pay (no causal benefit once in town), hence is not rescued.

FDT: commits to paying, because the computation ``would I pay?'' is
evaluated \emph{before} the rescue decision.  FDT agents get rescued and
pay, yielding higher utility than CDT agents who die in the desert.
\end{example}

\subsection{FDT and logical counterfactuals}

The central difficulty in FDT is defining the \emph{subjunctive
conditional} $P(e \mid \widehat{F}() = a)$.  This requires a theory of
\emph{logical counterfactuals}: what would the world be like if a
mathematical fact were different?  This connects directly to logical
induction (Section~\ref{sec:logical-induction}): one can use the logical
inductor's beliefs as the substrate for evaluating such counterfactuals.
Formalising this rigorously remains an open problem
(see Section~\ref{sec:open}).

%==========================================================================
\section{L\"ob's Theorem and the L\"obian Obstacle}\label{sec:lob}
%==========================================================================

\subsection{Statement}

\begin{theorem}[L\"ob's Theorem, 1955]\label{thm:lob}
Let $T$ be a recursively axiomatised theory extending $\PA$ and let
$\Prov_T(\ulcorner \cdot \urcorner)$ be its canonical provability
predicate.  For any sentence $\phi$:
\[
  T \vdash \bigl(\Prov_T(\ulcorner \phi \urcorner) \to \phi\bigr)
  \quad \Longrightarrow \quad
  T \vdash \phi.
\]
\end{theorem}

Equivalently: if $T$ proves ``if I can prove $\phi$, then $\phi$ is
true,'' then $T$ already proves $\phi$.  An agent that trusts its own
proofs can only do so for sentences it has already established.

\subsection{Proof sketch}

\begin{proof}[Proof sketch]
Assume $T \vdash \Prov_T(\ulcorner \phi \urcorner) \to \phi$.
\begin{enumerate}
  \item By the diagonal lemma, obtain a sentence $\psi$ such that
    $T \vdash \psi \leftrightarrow (\Prov_T(\ulcorner \psi \urcorner)
    \to \phi)$.
  \item Since $T \vdash \psi \to (\Prov_T(\ulcorner \psi \urcorner)
    \to \phi)$, and by the Hilbert--Bernays derivability conditions:
    \begin{itemize}
      \item $T \vdash \psi \implies T \vdash \Prov_T(\ulcorner \psi
        \urcorner)$ \hfill (D1)
      \item $T \vdash \Prov_T(\ulcorner \psi \to (\Prov_T(\ulcorner
        \psi \urcorner) \to \phi) \urcorner)$ \hfill (D1 on step 1)
      \item $T \vdash \Prov_T(\ulcorner \psi \urcorner) \to
        \Prov_T(\ulcorner \Prov_T(\ulcorner \psi \urcorner) \to \phi
        \urcorner)$ \hfill (D2: $\Prov$ distributes over $\to$)
      \item $T \vdash \Prov_T(\ulcorner \psi \urcorner) \to
        \Prov_T(\ulcorner \Prov_T(\ulcorner \psi \urcorner) \urcorner)
        $ \hfill (D3)
    \end{itemize}
  \item Combining these, $T \vdash \Prov_T(\ulcorner \psi \urcorner)
    \to \Prov_T(\ulcorner \phi \urcorner)$.
  \item By hypothesis, $T \vdash \Prov_T(\ulcorner \phi \urcorner) \to
    \phi$, so $T \vdash \Prov_T(\ulcorner \psi \urcorner) \to \phi$.
  \item By step 1's equivalence, $T \vdash \psi$.
  \item By D1, $T \vdash \Prov_T(\ulcorner \psi \urcorner)$.
  \item By step 4, $T \vdash \phi$.  \qedhere
\end{enumerate}
\end{proof}

\subsection{The L\"obian obstacle for cooperation}

Consider two agents, $A$ and $B$, each reasoning in $\PA$, who wish to
cooperate in a Prisoner's Dilemma.  Agent $A$ might try the strategy:
\begin{quote}
``Cooperate if I can prove that $B$ cooperates.''
\end{quote}
Symmetrically for $B$.  Let $a$ = ``$A$ cooperates'' and $b$ = ``$B$
cooperates.''  Then $A$'s strategy amounts to:
\[
  \PA \vdash \Prov_\PA(\ulcorner b \urcorner) \to a,
\]
and $B$'s to $\PA \vdash \Prov_\PA(\ulcorner a \urcorner) \to b$.

\begin{proposition}[Na\"ive mutual proof search fails]
The mutual proof search does not terminate: neither agent can establish
a proof of the other's cooperation without already having such a proof.
The circularity is exactly the structure L\"ob's theorem governs.
\end{proposition}

\subsection{The L\"obian handshake (Barasz et al.\ 2014)}

\begin{theorem}[Barasz, Christiano, Fallenstein, Herreshoff, LaVictoire, Yudkowsky~\cite{barasz2014robust}]
There exist agents $A$ and $B$ using the following strategy that
provably cooperate:

Agent $A$: ``Cooperate if and only if I can prove in $\PA + \Con(\PA)$
that $B$ cooperates.''

Agent $B$: ``Cooperate if and only if I can prove in $\PA$ that $A$
cooperates.''

\medskip\noindent
Here $A$ reasons in a \emph{stronger} system ($\PA + \Con(\PA)$) than
$B$ reasons about.  This asymmetry breaks the L\"obian cycle.
\end{theorem}

\begin{proof}[Proof sketch]
Working in $\PA + \Con(\PA)$:
\begin{enumerate}
  \item $A$ searches for a $\PA$-proof that $B$ cooperates.
  \item $B$'s code says: cooperate iff $\PA \vdash a$.
  \item Suppose for contradiction $\PA \vdash \neg a$.  Then $B$'s
    strategy is to defect, so $\PA \vdash \neg b$.  But $A$'s strategy
    in $\PA + \Con(\PA)$ would then also defect, so $\PA + \Con(\PA)
    \vdash \neg a$.  Since $\PA + \Con(\PA)$ is consistent (by
    $\Con(\PA)$), we have $\PA \nvdash a$, consistent with $\PA \vdash
    \neg a$.  But $A$'s code outputs defect, meaning $\PA + \Con(\PA)
    \vdash \neg a$ and $\PA \vdash \neg a$.  Now: $B$ defects, and $A$
    defects.  This is consistent, but we can show via the L\"obian
    structure that $\PA + \Con(\PA) \vdash b$, and hence $A$ cooperates,
    giving $\PA \vdash a$, so $B$ cooperates.
  \item The key insight: $A$'s use of the stronger system $\PA +
    \Con(\PA)$ allows it to verify that $\PA$ is consistent and hence
    that the mutual proof search terminates.
\end{enumerate}
\end{proof}

\begin{remark}
The L\"obian handshake demonstrates that \emph{hierarchies of logical
strength} can substitute for trust in multi-agent cooperation.  This is
a striking connection between proof theory and game theory.
\end{remark}

%==========================================================================
\section{Reflective Oracles}\label{sec:reflective-oracles}
%==========================================================================

\subsection{Motivation}

An embedded agent must reason about other agents who may themselves be
reasoning about it.  Reflective oracles, introduced by Fallenstein,
Taylor, and Christiano~\cite{fallenstein2015reflective}, provide a
mathematical framework in which agents can ``query'' a probabilistic
oracle about the output of machines that themselves have access to the
same oracle---without paradox.

\subsection{Definitions}

\begin{definition}[Oracle machine]
An \emph{oracle machine} $M^O$ is a Turing machine with access to an
oracle $O$.  We write $M^O(x)$ for the output of $M$ on input $x$
with oracle $O$.
\end{definition}

\begin{definition}[Reflective oracle]\label{def:reflective-oracle}
A \emph{reflective oracle} is a function
$O : \N \times \N \to \{0, 1\}$ such that for every pair of oracle
machines $(M_1, M_2)$ with index $e$ and input $x$:
\begin{enumerate}[label=(\roman*)]
  \item If $\Pr[M_1^O(x) \text{ halts and outputs } 1] > 1/2$,
    then $O(e, x) = 1$.
  \item If $\Pr[M_1^O(x) \text{ halts and outputs } 1] < 1/2$,
    then $O(e, x) = 0$.
  \item If the probability is exactly $1/2$, either answer is allowed.
\end{enumerate}
Here, the probability is over the internal randomness of $M_1^O$.  The
encoding $e$ packages the pair $(M_1, M_2)$ such that $O$ is queried
as: ``Does machine $M_1$ (which itself has access to $O$) output 1 with
probability $> 1/2$ on input $x$?''
\end{definition}

The self-referential nature is clear: $O$ answers questions about
machines that query $O$.

\subsection{Existence via Kakutani's fixed-point theorem}

\begin{theorem}[Existence of reflective oracles~\cite{fallenstein2015reflective}]
Reflective oracles exist.
\end{theorem}

\begin{proof}[Proof sketch]
Consider the space $\mathcal{X} = \{0,1\}^{\N \times \N}$ of all
possible oracles, equipped with the product topology.  This is a
compact, convex subset of a locally convex topological vector space
(via the embedding into $[0,1]^{\N \times \N}$).

Define the correspondence $F : \mathcal{X} \rightrightarrows \mathcal{X}$
by:
\[
  F(O) = \bigl\{ O' \in \mathcal{X} : O' \text{ is consistent with the
  probabilities induced by } O \bigr\}.
\]
More precisely, for each query $(e, x)$:
\begin{itemize}
  \item If $\Pr[M_e^O \text{ outputs } 1] > 1/2$, then $O'(e,x) = 1$.
  \item If $\Pr[M_e^O \text{ outputs } 1] < 1/2$, then $O'(e,x) = 0$.
  \item If equal to $1/2$, $O'(e,x)$ may be either $0$ or $1$.
\end{itemize}

One verifies:
\begin{enumerate}
  \item $F(O)$ is non-empty and convex for each $O$ (the ``tie-breaking''
    case provides the convexity).
  \item $F$ has closed graph (by continuity of probability measures
    under the product topology).
\end{enumerate}
By Kakutani's fixed-point theorem, $F$ has a fixed point: an oracle
$O^*$ with $O^* \in F(O^*)$.  This $O^*$ is a reflective oracle.
\end{proof}

\begin{remark}
Strictly, one works with a generalisation where the oracle returns
probabilities rather than bits, and the fixed-point is in a space of
probability measures. The argument above captures the essential idea.
\end{remark}

\subsection{Nash equilibrium with source code access}

\begin{theorem}[Nash equilibria under reflective oracles~\cite{fallenstein2015reflective}]
Consider a game where each player $i$ is an oracle machine $M_i^O$
choosing a (mixed) strategy.  Under any reflective oracle $O$, the
resulting strategy profile is a Nash equilibrium of the game.
\end{theorem}

\begin{proof}[Proof sketch]
Suppose player $i$ could improve by deviating to strategy $\sigma'$.
Encode $\sigma'$ as an oracle machine $M_i'^O$.  Then the reflective
oracle, being a fixed point, already accounts for the possibility that
$M_i^O$ outputs $\sigma'$---the oracle's answer for ``what does player
$i$ do?'' is consistent with $M_i^O$'s actual behaviour.  A profitable
deviation would break the fixed-point condition, contradiction.
\end{proof}

This is remarkable: agents with access to each other's source code
(and a reflective oracle) automatically reach Nash equilibria, without
any explicit equilibrium computation.

%==========================================================================
\section{Cartesian Frames}\label{sec:cartesian-frames}
%==========================================================================

\subsection{Motivation}

The classical agent--environment factorisation assumes a single, fixed
boundary.  Cartesian frames (Demski and Garrabrant~\cite{demski2020cartesian})
provide an algebraic framework for studying \emph{multiple} such
factorisations, their relationships, and when they are equivalent.

\subsection{Basic definitions}

\begin{definition}[Cartesian frame]
A \emph{Cartesian frame} is a triple $\mathcal{C} = (A, E, \boxtimes)$
where:
\begin{itemize}
  \item $A$ is a set (the \emph{agent}'s possible states/actions),
  \item $E$ is a set (the \emph{environment}'s possible states),
  \item $\boxtimes : A \times E \to W$ is a function to a fixed set of
    \emph{worlds} $W$.
\end{itemize}
We think of $\boxtimes(a, e) = w$ as: ``if the agent is in state $a$
and the environment is in state $e$, the resulting world is $w$.''
\end{definition}

\begin{definition}[Ensurables and preventables]
Given a Cartesian frame $(A, E, \boxtimes)$ and a set $S \subseteq W$:
\begin{itemize}
  \item $S$ is \emph{ensurable} if there exists $a \in A$ such that
    for all $e \in E$, $\boxtimes(a, e) \in S$.
  \item $S$ is \emph{preventable} if $W \setminus S$ is ensurable,
    i.e., there exists $a \in A$ such that for all $e \in E$,
    $\boxtimes(a, e) \notin S$.
  \item $S$ is \emph{controllable} if it is both ensurable and
    preventable.
\end{itemize}
\end{definition}

\begin{remark}
Ensurables capture what the agent can \emph{guarantee} regardless of
the environment.  They form a lattice under intersection (but not under
union in general).
\end{remark}

\begin{proposition}
The set of ensurables of a Cartesian frame $(A, E, \boxtimes)$ is
closed under arbitrary intersection: if $\{S_i\}_{i \in I}$ are all
ensurable, and $\bigcap_{i \in I} S_i \neq \emptyset$ in the range of
$\boxtimes$, then $\bigcap_{i \in I} S_i$ is ensurable.
\end{proposition}

\subsection{Morphisms and biextensional equivalence}

\begin{definition}[Morphism of Cartesian frames]
A \emph{morphism} from $(A_1, E_1, \boxtimes_1)$ to
$(A_2, E_2, \boxtimes_2)$ (over the same world set $W$) is a pair of
functions $f : A_1 \to A_2$ and $g : E_2 \to E_1$ such that for all
$a \in A_1$ and $e \in E_2$:
\[
  \boxtimes_2(f(a), e) = \boxtimes_1(a, g(e)).
\]
\end{definition}

Note the \emph{contravariance} in $E$: the environment map goes in the
opposite direction.  This reflects the adversarial/dual nature of the
environment relative to the agent.

\begin{definition}[Biextensional equivalence]
Two Cartesian frames are \emph{biextensionally equivalent} if they have
the same ensurables and preventables.  Equivalently, after quotienting
$A$ by the relation $a_1 \sim a_2 \iff \forall e.\, \boxtimes(a_1, e)
= \boxtimes(a_2, e)$ and quotienting $E$ by
$e_1 \sim e_2 \iff \forall a.\, \boxtimes(a, e_1) = \boxtimes(a, e_2)$,
the resulting frames are isomorphic.
\end{definition}

\begin{theorem}[Biextensional collapse]
Every Cartesian frame has a unique (up to isomorphism) biextensional
collapse: a biextensionally equivalent frame in which no two agent states
are observationally identical and no two environment states are
observationally identical.
\end{theorem}

\subsection{Sub-agents and composition}

\begin{definition}[Sub-agent]
A Cartesian frame $(A', E', \boxtimes')$ is a \emph{sub-agent} of
$(A, E, \boxtimes)$ if there is a morphism from the first to the
second that is injective on agent states.  Intuitively, $A'$ represents
a ``smaller'' agent embedded within $A$, with the remaining degrees of
freedom absorbed into the environment.
\end{definition}

\begin{definition}[Tensor product]
Given Cartesian frames $\mathcal{C}_1 = (A_1, E_1, \boxtimes_1)$ and
$\mathcal{C}_2 = (A_2, E_2, \boxtimes_2)$, their \emph{tensor product}
is:
\[
  \mathcal{C}_1 \otimes \mathcal{C}_2 =
    (A_1 \times A_2,\; E_1 \times E_2,\; \boxtimes)
\]
where $\boxtimes((a_1, a_2), (e_1, e_2)) =
  (\boxtimes_1(a_1, e_1), \boxtimes_2(a_2, e_2))$.
This models two agents acting independently in their respective
environments.
\end{definition}

\begin{remark}
The category of Cartesian frames (with morphisms as defined above) is
a $*$-autonomous category, connecting to linear logic.  The
multiplicative structure (tensor, par) captures composition and
decomposition of agency, while the additive structure captures choice.
\end{remark}

\subsection{Comparison with the dualistic model}

In the dualistic model of Section~\ref{sec:dualism}, the agent--environment
boundary is fixed.  In Cartesian frames, we can:
\begin{itemize}
  \item \emph{Refactor} the boundary: the same set of worlds may admit
    multiple non-isomorphic factorisations $(A_1, E_1)$ and $(A_2, E_2)$.
  \item \emph{Nest} agents: a sub-agent sits inside a larger agent,
    with the ``gap'' between them part of the environment.
  \item \emph{Compare} frames: biextensional equivalence tells us when
    two agent--environment decompositions are ``the same'' from the
    perspective of agency (same ensurables/preventables).
\end{itemize}

%==========================================================================
\section{Connections to Modern Agentic AI}\label{sec:connections}
%==========================================================================

The formalisms above were developed in the context of ``agent
foundations'' for superintelligence.  But they connect directly to
capabilities and failure modes of contemporary AI systems.

\subsection{LLM self-reference}

Large language models routinely produce text that references their own
capabilities, limitations, or outputs.  This creates G\"odelian
feedback loops:
\begin{itemize}
  \item \textbf{Self-evaluation prompts.}  ``Rate the quality of your
    previous answer'' asks the model to act as a fixed point of its own
    evaluation function.
  \item \textbf{Chain-of-thought faithfulness.}  When a model explains
    its reasoning, is the explanation causally related to the answer, or
    merely post-hoc?  This is a version of the embedded agency problem:
    the ``inner reasoner'' and ``outer reporter'' share the same
    weights.
  \item \textbf{L\"obian self-trust.}  A model asked ``Is your next
    output correct?'' faces a structure isomorphic to the L\"obian
    obstacle: it cannot verify its own correctness without already
    having the answer.
\end{itemize}

\subsection{Tool-using agents}

Modern LLM agents (e.g.\ ReAct, Toolformer, Claude with tool use) are
embedded in their environment: they emit actions (tool calls) that
modify state (files, databases, APIs) which then appears in their
context window.  This creates:
\begin{itemize}
  \item \textbf{Observation-action entanglement.}  The agent's output
    modifies the input it will receive, violating the dualistic
    assumption.
  \item \textbf{Logical uncertainty about tool outputs.}  Before calling
    a tool, the agent has uncertainty about the result that is partly
    logical (what will this code compute?) and partly empirical (what is
    the current database state?).
  \item \textbf{Multi-step planning under self-modification.}  An agent
    that writes code and then executes it is effectively modifying its
    own ``policy.''  The Cartesian frames formalism (Section~\ref{sec:cartesian-frames})
    offers a principled way to reason about which aspects of the world
    such an agent can control.
\end{itemize}

\subsection{Multi-agent systems}

When multiple LLM agents interact (e.g.\ in debate, negotiation, or
collaborative coding), the reflective oracle framework becomes directly
relevant:
\begin{itemize}
  \item Each agent's ``strategy'' depends on what other agents will do,
    which depends on what \emph{it} will do.
  \item Reflective oracles guarantee the existence of consistent
    fixed points (Nash equilibria) even with source-code transparency.
  \item The L\"obian handshake suggests that cooperation between
    formally reasoning agents requires \emph{asymmetric trust
    hierarchies}---a principle that could inform the design of
    multi-agent communication protocols.
\end{itemize}

\subsection{Constitutional AI and decision theory}

Constitutional AI (Bai et al.~\cite{bai2022constitutional}) trains
models to evaluate and revise their own outputs against a set of
principles.  This is a practical implementation of:
\begin{itemize}
  \item \textbf{Reflective stability}: the model's values should be
    a fixed point of its own revision process (cf.\ reflective oracles).
  \item \textbf{FDT-like reasoning}: the model reasons about ``what
    kind of output would a model following these principles produce?''
    rather than causally intervening on its own weights.
  \item \textbf{Logical induction over norms}: as the model processes
    more constitutional principles and examples, its ``normative
    beliefs'' should converge in a manner analogous to logical induction.
\end{itemize}

%==========================================================================
\section{Open Problems}\label{sec:open}
%==========================================================================

\begin{problem}[Logical counterfactuals]
Define a satisfactory theory of \emph{logical counterfactuals}: given a
mathematical structure $\mathcal{M}$ and a sentence $\phi$ true in
$\mathcal{M}$, what would $\mathcal{M}$ ``look like'' if $\phi$ were
false?  FDT requires such a theory but none of the existing proposals
(causal surgery on computation graphs, counterfactual mugging
responses, modal logic approaches) are fully satisfactory.  A formal
desideratum: the theory should be compatible with logical induction
and should handle nested counterfactuals
($P(e \mid \widehat{F}(\widehat{G}()) = a)$).
\end{problem}

\begin{problem}[Efficient logical induction]
Logical inductors as constructed in~\cite{garrabrant2016logical} are
computationally expensive (polynomial time per step, but with large
constants, and the polynomial degree grows with the number of traders).
Can one construct logical inductors that run in $O(n \log n)$ time per
step?  Alternatively, can neural networks be trained to approximate
logical inductors?
\end{problem}

\begin{problem}[Naturalized induction]
Extend Solomonoff induction to the case where the agent is part of the
environment.  A naturalized inductor would maintain a distribution over
``world-programs'' that include the agent itself as a subprocess.  The
challenge: the agent cannot enumerate world-programs longer than its own
description length, yet must reason about worlds more complex than
itself.  What are the right computability-theoretic or
complexity-theoretic constraints?
\end{problem}

\begin{problem}[Cartesian frames for neural networks]
The Cartesian frames formalism is set-theoretic.  Extend it to
continuous, differentiable settings appropriate for neural networks:
replace $A$ and $E$ with manifolds, $\boxtimes$ with a smooth map,
and characterise ``ensurables'' via reachability in the tangent space.
Can the resulting theory identify meaningful agent--environment
boundaries within a trained neural network?
\end{problem}

\begin{problem}[Multi-agent logical induction]
Extend logical induction to the multi-agent setting: multiple logical
inductors trading in a shared market, each with different computational
resources.  Does a multi-agent logical inductor converge to a correlated
equilibrium?  Can the L\"obian handshake be embedded in such a market?
\end{problem}

\begin{problem}[Reflective stability under self-improvement]
An agent that improves its own code creates a sequence of agents
$A_1, A_2, \ldots$.  Under what conditions does this sequence converge
to a ``reflectively stable'' agent $A_\infty$ that would not further
modify itself?  Reflective oracles guarantee existence of fixed points,
but do not address convergence of iterated self-improvement.  Can one
prove convergence under suitable smoothness or monotonicity conditions
on the improvement operator?
\end{problem}

\begin{problem}[Bounded L\"obian handshakes]
The L\"obian handshake requires reasoning in formal systems of
increasing logical strength.  In practice, agents have bounded
computational resources.  Can one achieve provable cooperation between
bounded agents (e.g.\ polynomial-time provers) without requiring the
hierarchy $\PA \subset \PA + \Con(\PA) \subset \cdots$?  Interactive
proof systems and probabilistically checkable proofs may be relevant.
\end{problem}

%==========================================================================
\section{Annotated Bibliography}\label{sec:bibliography}
%==========================================================================

\begin{enumerate}[label={[\arabic*]},leftmargin=2em]

\item \label{bib:demski2019} \textbf{Demski, A.\ and Garrabrant, S.}
  (2019). ``Embedded Agency.'' \textit{arXiv:1902.09469}.\\
  The foundational manifesto for the embedded agency research programme.
  Identifies four sub-problems (decision, logical uncertainty,
  high-bandwidth, subsystem alignment) and surveys the state of the art
  as of 2019.  Essential starting point.

\item \label{bib:garrabrant2016} \textbf{Garrabrant, S., Benson-Tilsen,
  T., Critch, A., Soares, N., and Taylor, J.} (2016). ``Logical
  Induction.'' \textit{arXiv:1609.03543}.\\
  Introduces the logical induction framework.  Proves existence of
  logical inductors and derives an extensive list of rationality
  properties (convergence, calibration, Gaifman condition,
  non-dogmatism, domination of CE sequences).  Technical and
  self-contained; the main theorem proof runs $\sim$60 pages.

\item \label{bib:soares2018} \textbf{Soares, N.\ and Yudkowsky, E.}
  (2018). ``Functional Decision Theory: A New Theory of Instrumental
  Rationality.'' \textit{arXiv:1710.05060}.\\
  Presents FDT and argues it outperforms CDT and EDT on Newcomb-like
  problems.  The philosophical framework is accessible; the formal
  details of logical counterfactuals are acknowledged as open.

\item \label{bib:barasz2014} \textbf{Barasz, M., Christiano, P.,
  Fallenstein, B., Herreshoff, M., LaVictoire, P., and Yudkowsky, E.}
  (2014). ``Robust Cooperation in the Prisoner's Dilemma: Program
  Equilibrium via Provability Logic.'' \textit{arXiv:1401.5577}.\\
  Introduces the L\"obian handshake: agents that cooperate by reasoning
  in formal systems of different strengths.  Elegant application of
  provability logic to game theory.

\item \label{bib:fallenstein2015} \textbf{Fallenstein, B., Taylor, J.,
  and Christiano, P.} (2015). ``Reflective Oracles: A Foundation for
  Game Theory in Artificial Intelligence.'' \textit{arXiv:1508.04145}.\\
  Defines reflective oracles and proves their existence via Kakutani's
  theorem.  Shows that agents with reflective oracle access automatically
  reach Nash equilibria.  Foundational for reasoning about interacting
  agents with mutual source-code access.

\item \label{bib:demski2020} \textbf{Demski, A.\ and Garrabrant, S.}
  (2020). ``Cartesian Frames.'' \textit{MIRI Technical Report}.\\
  Introduces Cartesian frames as a category-theoretic framework for
  agent--environment decompositions.  Defines ensurables, preventables,
  biextensional equivalence, and sub-agents.  Connects to $*$-autonomous
  categories and linear logic.

\item \label{bib:lob1955} \textbf{L\"ob, M.H.} (1955). ``Solution of a
  Problem of Leon Henkin.'' \textit{Journal of Symbolic Logic},
  20(2):115--118.\\
  The original proof of L\"ob's theorem.  Short and elegant.  Essential
  background for the L\"obian obstacle literature.

\item \label{bib:soares2015agent} \textbf{Soares, N.\ and Fallenstein,
  B.} (2017). ``Agent Foundations for Aligning Machine Intelligence with
  Human Interests.'' \textit{Machine Intelligence Research Institute
  Technical Report}.\\
  A high-level overview of the agent foundations research agenda, placing
  logical induction, decision theory, and reflective oracles in context.
  Good secondary source.

\item \label{bib:bai2022} \textbf{Bai, Y., Kadavath, S., Kundu, S.,
  et al.} (2022). ``Constitutional AI: Harmlessness from AI Feedback.''
  \textit{arXiv:2212.08073}.\\
  Introduces constitutional AI: training LLMs to critique and revise
  their own outputs.  Relevant as a practical instance of reflective
  stability in deployed systems.

\item \label{bib:critch2020} \textbf{Critch, A.\ and Krueger, D.}
  (2020). ``AI Research Considerations for Human Existential Safety
  (ARCHES).'' \textit{arXiv:2006.04948}.\\
  Surveys multi-agent safety, including formal models of boundaries
  between agents and environments.  Critch's boundaries formalism
  connects to Cartesian frames.

\item \label{bib:aaronson2016} \textbf{Aaronson, S.} (2016). ``The
  Complexity of Agreement.'' In \textit{Computational Complexity and
  Property Testing}, Springer.\\
  Discusses Aumann's agreement theorem from a computational complexity
  perspective.  Relevant to multi-agent logical induction: how quickly
  can bounded agents reach agreement on logical facts?

\item \label{bib:christiano2014} \textbf{Christiano, P.} (2014).
  ``Non-Omniscience, Probabilistic Inference, and Metamathematics.''
  \textit{arXiv:1408.2048}.\\
  An early approach to logical uncertainty predating logical induction.
  Uses reflection principles and probabilistic reasoning about
  arithmetic sentences.

\item \label{bib:yudkowsky2010} \textbf{Yudkowsky, E.} (2010).
  ``Timeless Decision Theory.'' \textit{MIRI Working Paper}.\\
  Precursor to FDT.  Introduces the idea of deciding based on
  ``what kind of algorithm am I'' rather than on causal or evidential
  considerations.  Less formally developed than FDT but historically
  important.

\item \label{bib:hutter2005} \textbf{Hutter, M.} (2005). \textit{Universal
  Artificial Intelligence: Sequential Decisions Based on Algorithmic
  Probability.}  Springer.\\
  The AIXI framework: optimal agent in a dualistic setting.  Essential
  reference for understanding what embedded agency is trying to move
  \emph{beyond}.  The mathematical formalism (Solomonoff induction +
  sequential decision theory) is the gold standard for the dualistic
  approach.

\item \label{bib:hadfield2017} \textbf{Hadfield-Menell, D., Russell,
  S., Abbeel, P., and Dragan, A.} (2017). ``The Off-Switch Game.''
  \textit{Proceedings of the AAAI Conference on AI}.\\
  Analyses the problem of an agent that might resist being shut off.
  Directly relevant to embedded agency: the off-switch is part of the
  environment the agent can influence, breaking the dualistic
  assumption.

\item \label{bib:taylor2016} \textbf{Taylor, J.} (2016). ``Reflective
  Alignment.'' \textit{MIRI Working Paper}.\\
  Extends reflective oracles toward alignment: can an agent verify
  that its successor (self-modified version) shares its values?  Uses
  logical induction and reflective oracles as building blocks.

\item \label{bib:gaifman2004} \textbf{Gaifman, H.} (2004). ``Reasoning
  with Limited Resources and Assigning Probabilities to Arithmetical
  Statements.'' \textit{Synthese}, 140(1--2):97--119.\\
  Early philosophical treatment of assigning probabilities to logical
  sentences under bounded resources.  Gaifman's condition (mentioned in
  logical induction properties) originates here.

\item \label{bib:kakutani1941} \textbf{Kakutani, S.} (1941). ``A
  Generalization of Brouwer's Fixed Point Theorem.'' \textit{Duke
  Mathematical Journal}, 8(3):457--459.\\
  The fixed-point theorem used to prove existence of reflective oracles.
  States that an upper hemicontinuous correspondence from a compact
  convex set to itself with non-empty convex values has a fixed point.

\end{enumerate}

%==========================================================================
% REFERENCES (for \cite commands)
%==========================================================================
\begin{thebibliography}{99}

\bibitem{demski2019embedded}
A.~Demski and S.~Garrabrant.
\newblock Embedded Agency.
\newblock \textit{arXiv:1902.09469}, 2019.

\bibitem{garrabrant2016logical}
S.~Garrabrant, T.~Benson-Tilsen, A.~Critch, N.~Soares, and J.~Taylor.
\newblock Logical Induction.
\newblock \textit{arXiv:1609.03543}, 2016.

\bibitem{soares2018fdt}
N.~Soares and E.~Yudkowsky.
\newblock Functional Decision Theory: A New Theory of Instrumental
  Rationality.
\newblock \textit{arXiv:1710.05060}, 2018.

\bibitem{barasz2014robust}
M.~Barasz, P.~Christiano, B.~Fallenstein, M.~Herreshoff, P.~LaVictoire, and
  E.~Yudkowsky.
\newblock Robust Cooperation in the Prisoner's Dilemma: Program Equilibrium
  via Provability Logic.
\newblock \textit{arXiv:1401.5577}, 2014.

\bibitem{fallenstein2015reflective}
B.~Fallenstein, J.~Taylor, and P.~Christiano.
\newblock Reflective Oracles: A Foundation for Game Theory in Artificial
  Intelligence.
\newblock \textit{arXiv:1508.04145}, 2015.

\bibitem{demski2020cartesian}
A.~Demski and S.~Garrabrant.
\newblock Cartesian Frames.
\newblock MIRI Technical Report, 2020.

\bibitem{bai2022constitutional}
Y.~Bai, S.~Kadavath, S.~Kundu, et~al.
\newblock Constitutional AI: Harmlessness from AI Feedback.
\newblock \textit{arXiv:2212.08073}, 2022.

\end{thebibliography}

\end{document}
